% AY2015 材料力学 〔II〕（転記）
% 出典: 04_材料力学/original/04_材料力学/2015/2015za.pdf
% 要約: 突き出しはり。左端A（壁に固定）からBまで L, BからC（自由端）まで L。
%       AB間に等分布荷重 f0, 自由端Cに集中荷重P（P < f0 L/4）。Bにローラー支持。
%       縦弾性係数E, 断面は幅b・高さhの長方形。支持反力・つり合い, SFD/BMD,
%       たわみ基礎式とたわみ角・たわみ, C点たわみを問う。

\section*{〔II〕}

図2に示す突き出しはりの自由端 C に集中荷重 $P$, AB 間に等分布荷重 $f_0$ が負荷される
場合について次の問い (1)\,$\sim$\,(3) に答えよ. 縦弾性係数を $E$, はりの断面形状を
幅 $b$, 高さ $h$ の長方形とする. ただし, 集中荷重は次の関係を持つとせよ. $(P<f_0L/4)$

\begin{center}
\begin{tikzpicture}[>=Latex]
  % 壁（左, 固定端A）
  \draw[thick] (0,-0.6) -- (0,1.8);
  \foreach \y in {-0.5,-0.2,...,1.75} {\draw (0,\y) -- (-0.28,\y+0.28);}
  % はり（A=0, B=4.5, C=9）
  \draw[line width=1.2pt] (0,0.1) -- (9,0.1);
  \draw[line width=1.2pt] (0,-0.1) -- (9,-0.1);
  % 等分布荷重 f0（AB間, 下向き）
  \draw (0,1.1) -- (4.5,1.1);
  \foreach \x in {0,0.5,...,4.5} {\draw[->] (\x,1.1) -- (\x,0.18);}
  \node at (2.25,1.4) {$f_0$};
  % ローラー支持 B
  \draw (4.5,-0.1) -- (4.2,-0.7) -- (4.8,-0.7) -- cycle;
  \draw (4.15,-0.85) -- (4.85,-0.85);
  % 集中荷重 P（C点 下向き）
  \draw[->,very thick] (9,1.5) -- (9,0.18);
  \node[right] at (9,1.3) {$P$};
  % 節点
  \node at (0.2,-0.45) {A};
  \node at (4.5,-1.05) {B};
  \node at (9,-0.45) {C};
  % 寸法 L, L
  \draw[<->] (0,-1.25) -- (4.5,-1.25) node[midway,fill=white]{$L$};
  \draw[<->] (4.5,-1.25) -- (9,-1.25) node[midway,fill=white]{$L$};
  % 座標 x
  \draw[->] (0,-1.7) -- (1.3,-1.7) node[right]{$x$};
\end{tikzpicture}

\vspace{1mm}
図2
\end{center}

\begin{enumerate}[label=(\arabic*)]
  \item 支持点 A, B に生じる支持力, 支持モーメントを図示し,
        力および曲げモーメントのつり合い式を求めよ.

  \item せん断力 $Q(x)$ と曲げモーメント $M(x)$ を求め,
        せん断力図（SFD）と曲げモーメント図（BMD）を示せ.

  \item たわみ $v(x)$ に関するはりの基礎式を示し, 積分によりたわみ角 $\dd v/\dd x$ と
        たわみ $v(x)$ を求め, C 点におけるたわみを求めよ.
\end{enumerate}
